\documentclass[oneside]{ncr_bundle}
\begin{document}

% ========================================================================
% Frontpage information 
% ========================================================================

% Editors of the conference proceedings as shown on the title page.
\author{Jonathan Archer \& Malcolm Reed (eds.)}

% Editors of the conference proceedings as to appear in the citation
\AuthorCitation{Archer, J., Reeds, M. (eds.)}

% Cover page (make your own)
\CoverPage{coverpages/coverpage.pdf}
% To credit the person who made the cover page, and/or credit photograph sources, use:
\CoverCredits{Cover attribution} 

\HeaderText{NCR DAYS 2018: The Future River. Deltares}


% Conference information
% ---------------------------------------------------------
\edition{2018}  % Edition, e.g. NCR DAYS 2022
\title{The future river}  % Title/theme of the conference
\date{8-9 February 2018}  % Date of the conference
\OrganisingPartner{1} % {1=Deltares, 2=Wageningen, 3=IHE Delft, 4=Twente.
%                        5=Utrecht, 6=rijkswaterstaat, 7=radboud, 8=TUD}
\PublicationNumber{42-2018}  % Request the number at the NCR Programme Secretary!


% Conference venue address & Contact
\ConferenceVenueAddress{
Deltares \\
Boussinesqweg 1\\
2629 HV Delft\\
P.O.\ 177\\
2600 MH Delft\\
The Netherlands
}

% Contact information. These are tab aligned. The first row *must* be \=
% subsequent rows \>, as shown below
\ConferenceVenueContact{
telephone: \=  +31\,88\,335\,82\,73\\
fax:       \>  +31\,88\,335\,85\,82\\
e-mail:    \>  info@deltares.nl\\
www:       \>  \url{https://www.deltares.nl}
}

% ========================================================================
% Frontpage generation, do not edit
% ========================================================================

% The \ncrtitle command generates the front page
\ncrtitle
       
% ========================================================================
% Introduction (information about the conference)
% ========================================================================
\chapter{Preface}
The NCR DAYS 2018 present the 20th time this conference is organised!
\blindtext

With riverine thoughts,\\
the LOC

\chapter{Conference details}
	\section{Organising partner}
	This partner is the best partner since 1648...
	\blindtext
	\section{Venues}
	Our campus was founded by the Visigoths and is widely considered ...
	\blindtext
	\section{Local organising committee (LOC)}
	The LOC, colloquially refereed to as `The Avengers' consists of ...
	\blindtext
 

 \chapter{Program}
 % insert program here. 

% ==============================================================
% Abstracts
%
% To add abstracts, use the \insertabstract command:
%
%    \insertabstract{section}{TITLE}{AUTHOR(S)}{PATH/TO/PDF}
%
%  for each abstract, use the above command and modify TITLE, 
%  AUTHOR(S) and PATH/TO/PDF. Abstracts are inserted as sections
%
%
%  To compartimentalise abstracts by session, you can either
%  use the \chapter command. Alternatively, if you have prepared
%  session covers, you can use the \insertsessioncover command:
%
%    \insertsessioncover{chapter}{SESSIONTITLE}{PATH/TO/PDF}{INT}
%
% where PATH/TO/PDF is the link to the session cover and INT
% is the pagenumber in the pdf to use. 
%
% See example below for the use of these two commands. 
% ==============================================================

% insert a session cover for session 1
\insertsessioncover{chapter}{Session 1}{coverpages/sessioncover.pdf}{1}

% insert abstracts
\insertabstract{section}{Archer, J., ...}{Title of abstract 1}{abstracts/01-latex.pdf}
\insertabstract{section}{Archer, J., ...}{Title of abstract 2}{abstracts/02-word.pdf}

% alternatively, you can start a section using chapter
\chapter{Session 2}
\insertabstract{section}{Archer, J., ...}{Title of abstract 1}{abstracts/01-latex.pdf}


% ========================================================================
% Outroduction (information about NCR)
% ========================================================================
\chapter{About NCR}
\begin{multicols}{2}

\section{Objective of NCR}
NCR was founded as a formal cooperation between several Dutch institutions. On October 4, 2012, the partners of the `Netherlands Centre for River studies' (NCR, Dutch: `Nederlands Centrum voor Rivierkunde (NCR)') formally renewed their cooperation within NCR.

The objective of NCR is:

\begin{quote}
Doelstelling van het NCR is een samenwerking tot stand te brengen tussen de belangrijkste kennisgebruikers en kennisontwikkelaars in Nederland op het gebied van rivieren met als uiteindelijk resultaat een versterking van het kennispotentieel, de profilering van de (inter)nationale positie van het Nederlands rivieronderzoek en het versterken van het onderwijs en het wetenschappelijk onderzoek aan de Nederlandse Universiteiten ten behoeve van een betere inrichting en beheer van de Nederlandse rivieren.
\end{quote}

This translates to:

\begin{quote}
The objective of the NCR is to establish cooperation between the major knowledge suppliers and knowledge users in Netherlands in the field of river studies, with the ultimate aim of reinforcing the knowledge potential, promoting the international position of Dutch river research and strengthen the education \& scientific research at Dutch universities, to better design and manage Dutch rivers.
\end{quote}

\section{Domains}
\label{section:domains}
NCR encompasses all disciplines relevant to river studies as practices by its institutional partners. They include:
\begin{itemize}
    \item Hydrodynamics
    \item Sediment transport and morphology
    \item Fluvial geomorphology and sedimentology
    \item River ecology, restoration and water quality
    \item Governance and spatial planning
    \item Modelling, serious gaming and digitial twins
\end{itemize}


\section{NCR Organisation}

NCR has four main bodies. The mandates for all except Young NCR (YNCR) are documented in the \emph{“Overeenkomst Nederlands Centrum voor Rivierkunde 2012”}.

\begin{itemize}
\item \textbf{Programme Secretary (Dutch: Programmasecretaris)} \\
Safeguards the continuity of NCR activities, secretary to SB and PC; monitors of agreed actions by the SB and PC; reports of NCR finances; management of all NCR communications.
\item \textbf{Supervisory Board (Dutch: Commissie van Toezicht)}\\ 
Supervises the implementation of the cooperation agreement, settles disputes and approves the scientific programme.
\item \textbf{Program Committee (Dutch: Programmacommissie)} \\
Determines the scientific programme, stimulates and initiates proposals for activities and integration of knowledge, ideas, experiences and results.
\item \textbf{Young NCR (Dutch: Jong NCR)} \\
Established in December 2020, YNCR strives to strengthen the network of young/early-career scientists within NCR. 
\end{itemize}

\section{Institutional partners}
The partners of NCR are: Rijkswaterstaat (RWS), Universiteit Twente (UT), Radboud Universiteit Nijmegen (RUN), Deltares, Universiteit Utrecht (UU), Technische Universiteit (TUD), Wageningen University \& Research (WUR), Leiden University Collega (LUC) and Vrije Universiteit Amsterdam (VU).

Each partner contributes in-kind and/or in-cash to realise the annual programme.  

\section{Accountability}
NCR activities are annually documented in an annual programme, annual report and annual budget. The annual programme and budget are approved by the Supervisory board. 

\end{multicols}



\begin{longtable}{@{}lll@{}}
\caption{The composition of NCR Boards in 2023. Persons marked with * are chair persons.}\\
% Please add the following required packages to your document preamble:
% \usepackage[table,xcdraw]{xcolor}
% If you use beamer only pass "xcolor=table" option, i.e. \documentclass[xcolor=table]{beamer}
% \usepackage[normalem]{ulem}
% \useunder{\uline}{\ul}{}

% Please add the following required packages to your document preamble:
% \usepackage[table,xcdraw]{xcolor}
% If you use beamer only pass "xcolor=table" option, i.e. \documentclass[xcolor=table]{beamer}
% \usepackage[normalem]{ulem}
% \useunder{\uline}{\ul}{}

\rowcolor[HTML]{629194} 
{\color[HTML]{FFFFFF} \textbf{Person}} & {\color[HTML]{FFFFFF} \textbf{Partner   institution}} & {\color[HTML]{FFFFFF} \textbf{Email}}                        \\
\endhead
\rowcolor[HTML]{C9D6D6} 
\textbf{Programme Secretary}           &                                                       &                                                              \\
\rowcolor[HTML]{B9C6D5} 
dr. ir. Koen Berends*                  & Deltares                                              & {\color[HTML]{0563C1} {\ul koen.berends@deltares.nl}}        \\
\rowcolor[HTML]{C9D6D6} 
ir. Anna Kosters                       & Deltares                                              & {\color[HTML]{0563C1} {\ul anna.kosters@deltares.nl}}        \\
\rowcolor[HTML]{B9C6D5} 
\textbf{Supervisory board}             &                                                       &                                                              \\
\rowcolor[HTML]{C9D6D6} 
prof. dr. Jaap Kwadijk*                & Deltares                                              & {\color[HTML]{0563C1} {\ul jaap.kwadijk@deltares.nl}}        \\
\rowcolor[HTML]{B9C6D5} 
prof. dr. Hans Middelkoop              & Utrecht University                                    & {\color[HTML]{0563C1} {\ul h.middelkoop@uu.nl}}              \\
\rowcolor[HTML]{C9D6D6} 
dr ir. Astrid Blom                     & Delft University of   Technology                      & {\color[HTML]{0563C1} {\ul astrid.blom@tudelft.nl}}          \\
\rowcolor[HTML]{B9C6D5} 
prof. dr. Ton  Hoitink                 & Wageningen University   \& Research                   & {\color[HTML]{0563C1} {\ul ton.hoitink@wur.nl}}              \\
\rowcolor[HTML]{C9D6D6} 
prof. dr. Suzanne Hulscher             & University of Twente                                  & {\color[HTML]{0563C1} {\ul s.j.m.h.hulscher@utwente.nl}}     \\
\rowcolor[HTML]{B9C6D5} 
drs. Monique Busnach-Blankers          & Rijkswaterstaat                                       & {\color[HTML]{0563C1} {\ul monique.busnach-blankers@rws.nl}} \\
\rowcolor[HTML]{C9D6D6} 
dr. Wilco Verberk                      & Radboud University   Nijmegen                         & {\color[HTML]{0563C1} {\ul w.verberk@science.ru.nl}}         \\
\rowcolor[HTML]{B9C6D5} 
\textit{Vacant}                        & Vrije Universiteit   Amsterdam                        & {\color[HTML]{0563C1} {\ul }}                                \\
\rowcolor[HTML]{C9D6D6} 
\textbf{Programme committee}           &                                                       &                                                              \\
\rowcolor[HTML]{B9C6D5} 
drs. Evelien van Eijsbergen*           & Rijkswaterstaat                                       & {\color[HTML]{0563C1} {\ul evelien.van.eijsbergen@rws.nl}}   \\
\rowcolor[HTML]{C9D6D6} 
dr. Susanne Quartel                   &                                                       & {\color[HTML]{0563C1} {\ul susanne.quartel@rws.nl}}          \\
\rowcolor[HTML]{B9C6D5} 
dr. Laura Stancanelli                  & Delft University of   Technology                      & {\color[HTML]{0563C1} {\ul L.M.Stancanelli@tudelft.nl}}      \\
\rowcolor[HTML]{C9D6D6} 
dr. Heleen Vreugdenhil                 &                                                       & {\color[HTML]{0563C1} {\ul heleen.vreugdenhil@deltares.nl}}  \\
\rowcolor[HTML]{B9C6D5} 
dr. Victor Chavarrias                  & Deltares                                              & {\color[HTML]{0563C1} {\ul victor.chavarrias@deltares.nl}}   \\
\rowcolor[HTML]{C9D6D6} 
drs. Valesca Harezlak                  &                                                       & {\color[HTML]{0563C1} {\ul valesca.harezlak@deltares.nl}}    \\
\rowcolor[HTML]{B9C6D5} 
dr. Frank Collas                       & Radboud University   Nijmegen                         & {\color[HTML]{0563C1} {\ul F.Collas@science.ru.nl}}          \\
\rowcolor[HTML]{C9D6D6} 
dr. Menno Straatsma                    & Utrecht University                                    & {\color[HTML]{0563C1} {\ul m.w.straatsma@uu.nl}}             \\
\rowcolor[HTML]{B9C6D5} 
dr. Jord Warmink                       & University of Twente                                  & {\color[HTML]{0563C1} {\ul j.j.warmink@utwente.nl}}          \\
\rowcolor[HTML]{C9D6D6} 
dr. Kryss Waldschläger                 & Wageningen   University \& Research                   & {\color[HTML]{0563C1} {\ul kryss.waldschlager@wur.nl}}       \\
\rowcolor[HTML]{B9C6D5} 
dr. Willem Toonen                      & Vrije Universiteit   Amsterdam                        & {\color[HTML]{0563C1} {\ul w.h.j.toonen@vu.nl}}              \\
\rowcolor[HTML]{C9D6D6} 
dr. Paolo Scussolini                   &                                                       & {\color[HTML]{0563C1} {\ul p.scussolini@vu.nl}}              \\
\rowcolor[HTML]{B9C6D5} 
dr. Paul Hudson                        & Leiden University   College                           & {\color[HTML]{0563C1} {\ul p.f.hudson@luc.leidenuniv.nl}}    \\
\rowcolor[HTML]{C9D6D6} 
\textbf{Young NCR}                     &                                                       &                                                              \\
\rowcolor[HTML]{B9C6D5} 
Niek van der Sleen MSc*                & Rijkswaterstaat-WVL                                   & {\color[HTML]{0563C1} {\ul niek.vander.sleen@rws.nl}}        \\
\rowcolor[HTML]{C9D6D6} 
ir. Anna Kosters                       & Deltares                                              & {\color[HTML]{0563C1} {\ul anna.kosters@deltares.nl}}        \\
\rowcolor[HTML]{B9C6D5} 
Jelle Dercksen MSc                     & Delft University of   Technology                      & {\color[HTML]{0563C1} {\ul J.A.Dercksen@tudelft.nl}}         \\
\rowcolor[HTML]{C9D6D6} 
Henry Amorocho Daza                    & IHE Delft          & {\color[HTML]{0563C1} {\ul }}                                \\
\rowcolor[HTML]{B9C6D5} 
Tatiana Edler MSc                      & Utrecht University                                    & {\color[HTML]{0563C1} {\ul t.edler@uu.nl}}                   \\
\rowcolor[HTML]{C9D6D6} 
ir. Fateme Ebrahimierami               & University of Twente                                  & {\color[HTML]{0563C1} {\ul f.ebrahimierami@utwente.nl}}      \\
\rowcolor[HTML]{B9C6D5} 
Iris Niesten MSc                       & Wageningen University   \& Research                   & {\color[HTML]{0563C1} {\ul iris.niesten@wur.nl}}             \\
\rowcolor[HTML]{C9D6D6} 
\textit{Vacant}                        & Radboud University   Nijmegen                         & {\color[HTML]{0563C1} {\ul }}                                \\
\rowcolor[HTML]{B9C6D5} 
\textit{Vacant}                        & Vrije Universiteit   Amsterdam                        & {\color[HTML]{0563C1} {\ul }}                                \\
\rowcolor[HTML]{C9D6D6} 
\textit{Vacant}                        & Leiden University   College                           & {\color[HTML]{0563C1} {\ul }}                                                   

\end{longtable}

\begin{figure}
\centering
\includegraphics[width=\textwidth]{figures/ncr_partners.png}
\end{figure}{}

% ========================================================================
% Notes (deprecated)
% Back when the proceedings where printed, the below command was used
% to generate 'notes' pages
% ========================================================================
%\notespages
	
\end{document}